\begin{figure}[tbp]
\centering
\begin{tikzpicture}[
  every node/.style={font=\footnotesize},
  live/.style={fill=allocbuddy!25},
  leak/.style={pattern=north east lines, pattern color=allocbuddy!80!black},
  cur/.style={-{Stealth[length=5pt]}, very thick},
  curc/.style={orange!85!black},
]
% One contiguous region, drawn growing upwards; each column is a snapshot.
\newcommand{\arU}{0.45}
\newcommand{\arW}{1.1}
\newcommand{\arpool}[1]{\draw[fill=gray!6] (#1,0) rectangle (#1+\arW,8*\arU);}
\newcommand{\arblock}[5]{% x, from, to, style, label
  \draw[#4] (#1,#2*\arU) rectangle (#1+\arW,#3*\arU);
  \node at (#1+\arW/2,{(#2+#3)/2*\arU}) {#5};}
\newcommand{\arcur}[2]{\draw[cur, curc] (#1+\arW+0.45,#2*\arU) -- (#1+\arW+0.02,#2*\arU);}
\newcommand{\arstep}[2]{\node[align=center, anchor=north, text width=2.3cm] at (#1+\arW/2,-0.12) {#2};}
\newcommand{\arA}{{\fboxsep=1pt\colorbox{white}{A}}}

% (a) empty
\arpool{0}
\arcur{0}{0}
\arstep{0}{(a) empty}

% (b) three allocations, bumped one after the other
\arpool{2.4}
\arblock{2.4}{0}{2.4}{live}{A}\arblock{2.4}{2.4}{3.4}{live}{B}\arblock{2.4}{3.4}{5.4}{live}{C}
\arcur{2.4}{5.4}
\arstep{2.4}{(b) allocate\\A, B, C}

% (c) free C: it is the newest, so the cursor moves back
\arpool{4.8}
\arblock{4.8}{0}{2.4}{live}{A}\arblock{4.8}{2.4}{3.4}{live}{B}
\arcur{4.8}{3.4}
\arstep{4.8}{(c) free C\\{\scriptsize\itshape cursor back}}

% (d) free A: not the newest, so nothing happens
\arpool{7.2}
\arblock{7.2}{0}{2.4}{leak}{\arA}\arblock{7.2}{2.4}{3.4}{live}{B}
\arcur{7.2}{3.4}
\arstep{7.2}{(d) free A\\{\scriptsize\itshape nothing happens}}

% (e) allocate D: smaller than A's old space, but still goes at the cursor
\arpool{9.6}
\arblock{9.6}{0}{2.4}{leak}{\arA}\arblock{9.6}{2.4}{3.4}{live}{B}\arblock{9.6}{3.4}{5.0}{live}{D}
\arcur{9.6}{5.0}
\arstep{9.6}{(e) allocate D\\{\scriptsize\itshape goes on top}}

% (f) deallocate_all
\arpool{12.0}
\arcur{12.0}{0}
\arstep{12.0}{(f) reset\\{\scriptsize\texttt{deallocate\_all}}}

% legend
\begin{scope}[shift={(2.1,-1.55)}]
  \draw[cur, curc] (0,0.125) -- (0.45,0.125); \node[right] at (0.5,0.125) {cursor};
  \draw[live] (2.0,0) rectangle (2.4,0.25); \node[right] at (2.45,0.125) {live};
  \draw[leak] (3.55,0) rectangle (3.95,0.25); \node[right] at (4.0,0.125) {freed, held until reset};
  \draw[fill=gray!6] (7.95,0) rectangle (8.35,0.25); \node[right] at (8.4,0.125) {free};
\end{scope}
\end{tikzpicture}
\caption{The arena allocator in IncludeOS, in six snapshots from left
  to right. The region is contiguous, and addresses grow upwards. A single
  \emph{cursor} (orange arrow) marks the first free byte, and allocating
  just moves it up by the requested size (b). Freeing only does something
  for the newest allocation, which moves the cursor back (c). Freeing
  anything older is accepted but does nothing, so A stays occupied until
  the whole arena is reset (d, f). A new allocation always goes at the
  cursor, even when it would fit in A's old space (e).}
\label{fig:arena-alloc}
\end{figure}
